\begin{ouverture}[geometrie]{12}{Pythagore et distances}{Partie 4 · Géométrie}
{Triangles rectangles · Réciproque · Distances · Modélisation}
\objectif{Calculer une longueur dans un triangle rectangle.}
\objectif{Prouver ou réfuter qu’un triangle est rectangle.}
\objectif{Construire un modèle de distance.}
\prerequis{Carrés, racines carrées et unités de longueur.}
\end{ouverture}
\pagestyle{otMain}\essentiel\fontsize{11}{15}\selectfont
\essSection{1}{Calculer une longueur}
\begin{propriete}[Théorème de Pythagore]Si ABC est rectangle en A, alors $BC^2=AB^2+AC^2$. BC, opposé à l’angle droit, est l’hypoténuse et le plus long côté.\end{propriete}
\begin{savoirfaire}[Chercher l’hypoténuse]Pour $AB=6$ cm et $AC=8$ cm, $BC^2=36+64=100$, donc $BC=\sqrt{100}=10$ cm. Une longueur est positive : on ne retient pas $-10$.\end{savoirfaire}
\begin{savoirfaire}[Chercher un côté de l’angle droit]Si $BC=13$ cm et $AB=5$ cm, $AC^2=BC^2-AB^2=169-25=144$, d’où $AC=12$ cm.\end{savoirfaire}
\begin{attention}Le théorème exige un angle droit connu. Une apparence de dessin ne suffit pas. Lorsqu’une racine n’est pas entière, conserver la valeur exacte puis arrondir à la précision demandée.\end{attention}
\clearpage
\essSection{2}{Décider et modéliser}
\begin{propriete}[Réciproque]Dans un triangle, si le carré du plus grand côté égale la somme des carrés des deux autres, alors le triangle est rectangle. L’angle droit est opposé au plus grand côté.\end{propriete}
\begin{propriete}[Contraposée]Si cette égalité n’est pas vérifiée, le triangle n’est pas rectangle. Écrire séparément les deux calculs permet une comparaison claire.\end{propriete}
\begin{savoirfaire}[Une distance dans un repère orthonormé]Entre deux points, construire un triangle rectangle avec un segment horizontal et un segment vertical. Leurs longueurs sont les écarts d’abscisses et d’ordonnées. Appliquer Pythagore. Un repère non orthonormé ne permet pas ce calcul direct.\end{savoirfaire}
\begin{savoirfaire}[Un problème concret]Identifier les points, préciser les verticales et horizontales, dessiner l’angle droit puis traduire les dimensions. Vérifier que la réponse correspond bien à la longueur demandée et reste cohérente avec les autres côtés.\end{savoirfaire}
\clearpage\rubrique{Découverte}
\begin{blocLivre}{ex-12-1}\begin{exercice}[1]{Reconnaître l’hypoténuse}
ABC est rectangle en A. Nommer l’hypoténuse et écrire l’égalité de Pythagore.
\label{fig-ex-12-1}\figureGeo{pyth-01}
\end{exercice}\end{blocLivre}
\begin{blocLivre}{ex-12-2}\begin{exercice}[1]{Un triangle classique}
ABC est rectangle en A, avec AB = 6 cm et AC = 8 cm. Calculer BC.
\label{fig-ex-12-2}\figureGeo{pyth-02}
\end{exercice}\end{blocLivre}
\begin{blocLivre}{ex-12-3}\begin{exercice}[1]{Le côté manquant}
ABC est rectangle en A. BC = 13 cm et AB = 5 cm. Calculer AC.
\label{fig-ex-12-3}\figureGeo{pyth-03}
\end{exercice}\end{blocLivre}
\begin{blocLivre}{ex-12-4}\begin{exercice}[1]{Garder une racine}
ABC est rectangle en A. AB = 4 cm, AC = 7 cm. Donner BC exactement puis au dixième.
\label{fig-ex-12-4}\figureGeo{pyth-04}
\end{exercice}\end{blocLivre}
\begin{blocLivre}{ex-12-5}\begin{exercice}[1]{Tester la réciproque}
Dans ABC, AB = 9 cm, AC = 12 cm et BC = 15 cm. Le triangle est-il rectangle ?
\label{fig-ex-12-5}\figureGeo{pyth-05}
\end{exercice}\end{blocLivre}

\clearpage\rubrique{Entraînement}
\begin{blocLivre}{ex-12-6}\begin{exercice}[2]{Refuser une fausse égalité}
AB = 5 cm, AC = 7 cm et BC = 9 cm. ABC est-il rectangle ?
\label{fig-ex-12-6}\figureGeo{pyth-06}
\end{exercice}\end{blocLivre}
\begin{blocLivre}{ex-12-7}\begin{exercice}[2]{Une échelle contre un mur}
Le mur AC est vertical et le sol AB horizontal. Une échelle BC de 5 m est posée à 1,4 m du mur. Quelle hauteur AC atteint-elle ?
\label{fig-ex-12-7}\figureGeo{pyth-07}
\end{exercice}\end{blocLivre}
\begin{blocLivre}{ex-12-8}\begin{exercice}[2]{Diagonale d’un rectangle}
Les côtés perpendiculaires d’un rectangle mesurent 15 cm et 20 cm. ABC en représente une moitié. Calculer la diagonale BC et le périmètre du rectangle.
\label{fig-ex-12-8}\figureGeo{pyth-08}
\end{exercice}\end{blocLivre}
\begin{blocLivre}{ex-12-9}\begin{exercice}[2]{Deux points du plan}
Dans un repère orthonormé d’unité 1 cm, A(1;2) et B(7;10). Calculer AB.
\label{fig-ex-12-9}\figureGeo{pyth-repere}
\end{exercice}\end{blocLivre}
\begin{blocLivre}{ex-12-10}\begin{exercice}[2]{Un angle à identifier}
AB = 10 cm, AC = 6 cm et BC = 8 cm. Prouver que le triangle est rectangle et préciser où.
\label{fig-ex-12-10}\figureGeo{pyth-10}
\end{exercice}\end{blocLivre}

\clearpage\rubrique{Problèmes}
\begin{blocLivre}{ex-12-11}\begin{exercice}[3]{Une voile et son aire}
Une voile ABC est rectangle en A. AB = 2,4 m et BC = 4 m. Calculer AC puis son aire.
\label{fig-ex-12-11}\figureGeo{pyth-11}
\end{exercice}\end{blocLivre}
\begin{blocLivre}{ex-12-12}\begin{exercice}[3]{Passer un panneau}
Une ouverture rectangulaire a pour côtés 0,9 m et 1,2 m. Sa diagonale BC est-elle assez longue pour une tige rigide de 1,48 m, d’épaisseur négligeable, placée dans le plan de l’ouverture ?
\label{fig-ex-12-12}\figureGeo{pyth-12}
\end{exercice}\end{blocLivre}
\begin{blocLivre}{ex-12-13}\begin{exercice}[3]{Une rampe}
Une rampe BC monte de AC = 0,75 m sur une avancée horizontale AB = 10 m. Calculer la longueur de rampe au centimètre. La pente est définie ici par $100\,AC/AB$. Calculer cette pente.
\label{fig-ex-12-13}\figureGeo{pyth-13}
\end{exercice}\end{blocLivre}
\begin{blocLivre}{ex-12-14}\begin{exercice}[3]{Deux chemins}
Un parc rectangulaire mesure 120 m sur 160 m. On compare le trajet suivant deux côtés adjacents au trajet diagonal. Quelle distance économise-t-on ?
\label{fig-ex-12-14}\figureGeo{pyth-parc}
\end{exercice}\end{blocLivre}
\begin{blocLivre}{ex-12-15}\begin{exercice}[3]{Une hauteur de triangle}
ABC est isocèle en C : CA = CB = 13 cm et AB = 10 cm. H est le milieu de AB. Calculer CH puis l’aire de ABC.
\label{fig-ex-12-15}\figureGeo{pyth-iso}
\end{exercice}\end{blocLivre}

\clearpage
